\documentclass[./测控三级项目.tex]{subfiles} % 使用 subfiles 文档类，指定主文档

\begin{document}
	\subsection{题目重述}
		\begin{compactenum}
			\item 分别设计符合要求的滤波器电路，用Multisim对滤波器的频率特性进行仿真。
			\item 搭建实测电路，分别为上述实测电路引入一个20mV正弦信号，调整信号频率，在0.1Hz-1MHz 范围内测试30个以上关键频率点的输出波形，记录输出信号幅值变化，用MATLAB描点绘制Bode图，用术语描述滤波器的频率特性（使用对数坐标，纵坐标用dB刻度，取点时可以是不等间距的）。分析设计电路是否完成功能，并进行性能分析。
			\item 做一阶、二阶、三阶、四阶有源低通滤波器的对比Bode图，叙述差异。
		\end{compactenum}
	\subsection{题目分析与实施方案}
		\subsubsection{第一问}
			第一问要求使用四阶低通滤波器抑制题6[1]e)中的中高频信号。将四阶低通滤波器看作一个一阶低通有源滤波器和一个Sallen-Key结构二阶有源低通滤波器。其中由于需要保留的低频信号频率为200Hz较低，故不采用同相放大，而是采用电压跟随器。\par
			电路确定后，通过计算有源低通网络的传递函数，得到对应相频特性函数。将需要截止频率代入特性曲线的截止频率，便可以推算出对应电路参数，并根据实际情况，选用相关精度的标准阻值电阻。 \par
			然后，使用 Multisim 软件仿真对应电路。
		\subsubsection{第二问}
			根据第一问搭建的电路，向四阶低通滤波器等对数频率，分别测试输出信号的幅值变化三十组。\par 
			使用 MATLAB 的 plot 函数，描点绘制出测试得到的参数，以得到相频伯德图曲线。 \par
			根据绘制的伯德图曲线，分析该四阶低通滤波器的频率特性。
		\subsubsection{第三问}
			根据上述题目做出的一阶、二阶、三阶、四阶有源低通滤波器的仿真图像，综合观察伯德图图像得到结论，叙述差异
	\subsection{理论分析与仿真}
		\subsubsection{第一问}
			图\ref{fig:四阶低通滤波器}是四阶低通滤波器的原理图。
			\begin{figure}[h]
				\centering
				\includegraphics[width=0.45\textwidth]{四阶低通滤波器原理图}
				\caption{四阶低通滤波器原理图}
				\label{fig:四阶低通滤波器}
			\end{figure}\par
			四阶低通滤波器由一个一阶低通有源滤波器和一个Sallen-Key结构二阶有源低通滤波器级联而成。\par
			式\ref{eq:Sallen-Key结构二阶有源低通滤波器传递函数}表达了Sallen-Key结构二阶有源低通滤波器的传递函数。因此将其开平方便可以得到对应的传递函数：
			\begin{gather}
				G(s)=[(1+\frac{R_5}{R_6})\frac{1}{R_1C_1R_2C_2s^2+(R_1C_2+R_2C_2-\frac{R_5}{R_6}R_1C_1)s+1}]^2
			\end{gather}\par
			为了方便计算，将其传递函数分为两个基本部分，进行运算。\par
			Sallen-Key结构二阶有源低通滤波器的传递函数分解如下：
			\begin{align}
				G_1(s)=\frac{v_{o1}}{v_{i1}}=(1+\frac{R_5}{R_6})\frac{1}{R_1C_1R_2C_2s^2+(R_1C_2+R_2C_2-\frac{R_5}{R_6}R_1C_1)s+1}
			\end{align}\par
			根据式\ref{eq:Sallen-Key结构二阶有源低通滤波器截止频率}可给出对应归一化参数公式为:
			\begin{gather}
				\begin{dcases}
					H_0=1+\frac{R_5}{R_4}	\\
					\omega_0 =\frac{1}{\sqrt{R_2C_2R_3C_3}}		\\
					Q =	\frac{R_5}{2R_5-R_4}
				\end{dcases}
				\label{eq:Sallen-Key结构二阶有源低通滤波器截止频率2}
			\end{gather}\par
			以上理论完备，在Multisim软件中进行仿真实验。\par
			由于需要滤除的中频信号为500Hz，需要保留的低频信号为200Hz。所以，设置截止频率为$f_0=350$Hz。\par
			设置电容容值为$C_1=C_2=C_3=C_4=100$nF(Standeard Value,$5\%$)\par
			根据式\ref{eq:Sallen-Key结构二阶有源低通滤波器截止频率2}，电阻阻值$R_1=R_2=R_3=R_4=4547.284088\Omega$，与标准电阻值表比对后，实际选用电阻阻值为$R_1=R_2=R_3=R_4=4.53$k$\Omega$(Standeard Value,$1\%$)的电阻。\par
			同时，我们希望使用同相放大器对低频信号进行补偿，即当$f_L=200$Hz时，增益为1：
			\begin{gather}
				A(\omega_M)=\lvert G(j\omega_M) \rvert =1
			\end{gather}\par
			将相关参数代入计算，整理后得到：
			\begin{gather}
				R_5=R_7=0.2186734297R_6=0.2186734297R_8
			\end{gather}\par
			设置电阻$R_6$、$R_8$的阻值为$R_6=R_8=150$k$\Omega$(Standeard Value,$1\%$)，则$R_5$、$R_7$阻值应为$R_5=R_7=32.80101445$k$\Omega$，与标准电阻值表比对后，实际选用电阻阻值为$R_5=R_7=33.2$k$\Omega$(Standeard Value,$1\%$)的电阻。\par
			将要求的信号接入该电路，得到仿真电路图，如图\ref{fig:四阶低通滤波器仿真原理图}所示。
			\begin{figure}[H]
				\centering
				\includegraphics[width=0.6\textwidth]{./18.第十八题示波器现象/18仿真原理图}
				\caption{四阶低通滤波器仿真原理图}
				\label{fig:四阶低通滤波器仿真原理图}
			\end{figure}\par
			对输入和输出信号做傅里叶变换，得到对应频谱曲线，如图\ref{fig:仿真四阶低通滤波器频谱曲线}所示。
			\begin{figure}[h]
				\centering
				\includegraphics[width=0.6\textwidth]{./18.第十八题示波器现象/18仿真频率曲线}
				\caption{仿真四阶低通滤波器频谱曲线}
				\label{fig:仿真四阶低通滤波器频谱曲线}
			\end{figure}\par
			对于低频信号（200Hz正弦信号），幅值从0.399985降到了0.402663，可求得增益为$A=20\lg(\frac{0.402663}{0.399985})=0.05793868248$dB；对于中频信号（500Hz正弦信号），幅值从0.399949降到了0.0799284，可求得增益为$A=20\lg(\frac{0.0799284}{0.399949})=-13.98628726$dB；对于高频信号（1kHz正弦信号），幅值从0.0999459降到了0.00195507，可求得增益为$A=20\lg(\frac{0.00195507}{0.0999459})=-34017205343$dB。\par
			这说明图\ref{fig:四阶低通滤波器}的四阶低通滤波器可以对高频信号和中频信号进行滤波操作。
	\subsection{实际电路测试结果与数据分析}
		\subsubsection{第二问}
			按题目要求引入幅值为20mV的正弦信号，并调整信号频率，在0.1Hz-1MHz区间测试30个以上关键频率点，测试结果如表\ref{tab:不同频率下四阶低通滤波器增益表}所示。
			\begin{table}[h!]
				\centering
				\caption{不同频率下四阶低通滤波器增益表}
				\label{tab:不同频率下四阶低通滤波器增益表}
				\begin{tabular}{|c|c|c|c|c|c|c|c|}
					\hline
					频率（Hz） & 增益  & 频率（Hz） & 增益 & 频率（Hz） & 增益 & 频率（Hz） & 增益\\
					\hline
					0.1&30.52&10&30.41&1000&0.3&100000&0.01\\\hline
					0.3&30.52&30&29.58&3000&0.01&300000&0.01\\\hline
					0.5&30.52&50&28.04&5000&0.01&500000&0.01\\\hline
					0.7&30.52&70&26&7000&0.01&700000&0.01\\\hline
					0.9&30.52&90&23.68&9000&0.01&900000&0.01\\\hline
					1&30.52&100&22.47&10000&0.01&1000000&0.01\\\hline
					3&30.51&300&6.5&30000&0.01&&\\\hline
					5&30.49&500&2.2&50000&0.01&&\\\hline
					7&30.47&700&0.9&70000&0.01&&\\\hline
					9&30.43&900&0.43&90000&0.01&&\\\hline
				\end{tabular}
			\end{table}	\par
			将表\ref{tab:不同频率下四阶低通滤波器增益表}中的数据输入MATLAB中进行Bode图生成，其图像如下图\ref{fig:四阶低通滤波器测试伯德图}所示。
			\begin{figure}[H]
				\centering
				\includegraphics[width=0.65\textwidth]{./18.第十八题示波器现象/四阶低通滤波器测试伯德图}
				\caption{测试生成四阶低通滤波器伯德图}
				\label{fig:四阶低通滤波器测试伯德图}
			\end{figure}\par
			根据实际选取电阻电容值参数，可以直接得到对应传递函数。利用MATLAB软件，得到该传递函数的伯德图如图\ref{fig:四阶低通滤波器实际伯德图}所示。
			\begin{figure}[H]
				\centering
				\includegraphics[width=0.65\textwidth]{./18.第十八题示波器现象/四阶低通滤波器实际伯德图}
				\caption{理论生成四阶低通滤波器伯德图}
				\label{fig:四阶低通滤波器实际伯德图}
			\end{figure}\par
			对比图\ref{fig:四阶低通滤波器测试伯德图}和图\ref{fig:四阶低通滤波器实际伯德图}中的相频曲线，可以发现除了曲线前者不如后者光滑外，其余基本一致，这说明实际测试结果符合理论推断。\par
			同时也注意到，在90Hz-3kHz范围内，增益大幅变化，从35下降到0的位置，这说明四阶低通滤波器具有通过低频信号的特点。
		\subsubsection{第三问}
			对比图\ref{fig:同相一阶有源低通滤波器测试伯德图}的同相一阶有源低通滤波器伯德图、图\ref{fig:Sallen-Key结构二阶有源低通滤波器测试伯德图}的Sallen-Key结构二阶有源低通滤波器伯德图、图\ref{fig:三阶低通滤波器测试伯德图}的三阶低通滤波器伯德图和图\ref{fig:四阶低通滤波器测试伯德图}的四阶低通滤波器伯德图可以发现，随着阶次的上升，滤波器下降速率增大，过渡带的带宽减小。
	\subsection{结论与收获}
	\textbf{游寅龙}：\par
	理论分析层面:查找了四阶低通滤波器的电路原理图，深入理解了四阶低通滤波器的电路原理。\par
	动手实践方面:能熟练使用信号发生器产生20mV正弦信号，并能运用示波器对结果进行观测和记录，提高了动手能力。\par
	精神层面: 数据的测量与记录使得耐心和毅力大大提升，让我处理事情更加沉稳。\par
	\textbf{郭嘉媛}:\par
	理论分析层面:了解其基本的电路组成和工作原理，能通过推导和分析传递函数，理解其在不同频率下的增益、相移等特性变化规律，增强了总结能力。\par
	动手实践层面:完成了Multsimn 仿真的搭建，能够根据理论分析实时调整搭建的电路，增强了独立解决问题的能力并丰富了解决问题的手段。\par
	精神层面:仿真的成功搭建，给人强烈的成就感，让我对测控电路的兴趣更深厚。\par
	\textbf{郑嘉宁}:\par
	理论分析层面:理解了滤波器截止频率在电路设计中的关键作用，能够根据不同的题目要求计算相应的截止频率，知识掌握更加熟练。\par
	动手实践层面:完成了对实测电路的搭建，能够根据要求调试电路，焊接能力有很大提升。\par
	精神层面:实测电路的搭建使我的抗压能力有了较大提升，让我能以积极的心态去面对问题、解决问题。\par
	
	
	
	
\end{document}